\begin{table}[htbp]
\centering
\small
\caption{The liquidity validation tests run on the five path and formula variants of Silicon Valley Bank: the light (G green, A amber, R red, n/e not evaluated, a dash where the test does not apply to the variant) and, in parentheses, the statistic; the run-off adequacy row is the ratio of modelled to observed and expected cumulative outflow at the two-day event horizon against the 85 per cent of the board review, which includes the outflow that management expected for the second day.}
\label{tab:alm-c16-battery}
\fitwidth{%
\begin{tabular}{lccccc}
\toprule
Test & \shortstack{LCR standard\\run-off} & LiST Adverse & LiST Extreme & 2023 fast run & \shortstack{HTM losses in\\the CBC} \\
\midrule
Compounding example, closed form at six decimals (absolute deviation) & \relax G ($3.66\times10^{-7}$) & \relax G ($3.66\times10^{-7}$) & \relax G ($3.66\times10^{-7}$) & \relax G ($3.66\times10^{-7}$) & \relax G ($3.66\times10^{-7}$) \\
Nine haircut cases (largest deviation) & \relax G (0) & \relax G (0) & \relax G (0) & \relax G (0) & \relax G (0) \\
Net liquidity position identity (residual) & -- & \relax G (0) & \relax G (0) & \relax G (0) & \relax G (0) \\
Survival monotone in severity (largest increase, days) & -- & -- & \relax G (0) & -- & -- \\
Cliff effect (share of assets) & -- & \relax A ($-$0.0222) & \relax A ($-$0.0298) & \relax n/e & \relax A ($-$0.0222) \\
Modelled over observed and expected outflow, board share & \relax A (0.506) & \relax R (0.108) & \relax R (0.115) & \relax R (0.201) & \relax R (0.108) \\
Held-to-maturity loss lowers the capacity (USD thousands) & -- & -- & -- & -- & \relax G ($1.45\times10^{7}$) \\
Revaluation against reported fair value & -- & -- & -- & -- & \relax R (0.282) \\
Run-off against generator truth (largest z) & -- & -- & -- & \relax G (0.537) & -- \\
Survival monotone in uninsured share (days) & -- & \relax G (0) & -- & -- & -- \\
\bottomrule
\end{tabular}}
\end{table}
